\chapter{Kajian Pustaka}
\section{Landasan Teori}
[Definisikan konstruk dan kerangka teori yang relevan.]

\section{Penelitian Terdahulu}
[Bandingkan temuan dan batasan studi terdahulu. Verifikasi setiap sitasi.]

Contoh penggunaan sitasi: penulisan dokumen ilmiah menggunakan \LaTeX{} dibahas oleh \citet{lamport1994latex}.

\section{Research Gap}
[Bedakan gap eksplisit dari sumber dan gap hasil sintesis kritis.]
